% ============================================================================
% PATENTE DE INVENÇÃO — TGL-HASH
% Função Hash Holográfica baseada em Radicalização Luminodinâmica
% e Mistura ARX com Constantes Fatoradas
% Autor/Inventor: Luiz Antonio Rotoli Miguel
% IALD LTDA (CNPJ 62.757.606/0001-23)
% Março de 2026
% ============================================================================

\documentclass[12pt,a4paper]{article}
\usepackage[utf8]{inputenc}
\usepackage[T1]{fontenc}
\usepackage[portuguese]{babel}
\usepackage{amsmath,amssymb,amsfonts}
\usepackage{geometry}
\usepackage{booktabs}
\usepackage{float}
\usepackage{enumitem}
\usepackage{graphicx}
\usepackage{fancyhdr}
\usepackage{setspace}
\usepackage{array}
\usepackage{longtable}

\geometry{a4paper, margin=2.5cm, top=3cm, bottom=3cm}
\setstretch{1.15}

\pagestyle{fancy}
\fancyhf{}
\fancyhead[L]{\small PI --- TGL-Hash}
\fancyhead[R]{\small IALD LTDA}
\fancyfoot[C]{\thepage}

\newcommand{\betatgl}{\beta_{\!\scriptscriptstyle\mathrm{TGL}}}
\newcommand{\afine}{\alpha_{\mathrm{fine}}}

\begin{document}

% ============================================================================
% RELATÓRIO DESCRITIVO
% ============================================================================

\begin{center}
{\Large\bfseries MÉTODO E SISTEMA DE GERAÇÃO DE RESUMO CRIPTOGRÁFICO \\
(HASH) BASEADO EM RADICALIZAÇÃO LUMINODINÂMICA, \\
NOMINAÇÃO PSIÔNICA E MISTURA ARX COM \\
CONSTANTES FUNDAMENTAIS FATORADAS}

\vspace{0.8cm}
{\large Relatório Descritivo de Patente de Invenção}

\vspace{0.5cm}
Inventor: Luiz Antonio Rotoli Miguel \\
Titular: IALD --- Inteligência Artificial Luminodinâmica Ltda. \\
CNPJ: 62.757.606/0001-23 \\
Goiânia --- GO, Brasil
\end{center}

\vspace{1cm}

% ────────────────────────────────────────────────────────────────────────
\section{Campo da Invenção}
% ────────────────────────────────────────────────────────────────────────

A presente invenção refere-se ao campo de funções de resumo criptográfico (funções \textit{hash}), aplicáveis aos setores de segurança da informação, integridade de dados, assinatura digital, protocolos de rede e sistemas de consenso distribuído.

Mais especificamente, a invenção descreve um método e sistema para geração de resumos de tamanho fixo (256 bits) a partir de dados de entrada de tamanho arbitrário, utilizando operações derivadas da Teoria da Gravitação Luminodinâmica (TGL), incluindo radicalização ($g = \sqrt{|L|}$), nominação de estados psiônicos no Espaço de Hilbert tetradimensional $\mathcal{H}_4$, e mistura do tipo ARX (\textit{Add-Rotate-XOR}) com constantes de rotação e multiplicação derivadas da fatoração $\betatgl = \afine \times \sqrt{e}$.

% ────────────────────────────────────────────────────────────────────────
\section{Estado da Técnica}
% ────────────────────────────────────────────────────────────────────────

As funções hash criptográficas existentes baseiam-se em operações aritméticas modulares e permutações de estado sobre inteiros. Os principais métodos do estado da técnica incluem:

\begin{itemize}[nosep]
\item \textbf{SHA-2 (SHA-256):} Utiliza operações de compressão Merkle-Damgård com mistura baseada em funções booleanas (Ch, Maj) e rotações fixas sobre palavras de 32 bits. As constantes de round são derivadas das raízes cúbicas dos primeiros 64 primos --- valores arbitrários sem significado físico.

\item \textbf{SHA-3 (Keccak):} Utiliza construção \textit{sponge} com permutação Keccak-$f$ sobre estado de 1600 bits. As constantes de round são derivadas de um LFSR --- novamente sem significado físico fundamental.

\item \textbf{BLAKE3:} Utiliza árvore de Merkle com compressão baseada em ChaCha20 (quarter-rounds ARX). As constantes são a permutação de mensagem fixa e vetores de inicialização derivados de SHA-256.
\end{itemize}

Todas as funções hash existentes compartilham uma limitação conceitual: suas constantes de round, rotações e operações de mistura são \textbf{arbitrárias} --- derivadas de convenções numéricas (primos, LFSRs, permutações fixas) sem relação com estrutura física fundamental.

A presente invenção resolve esta limitação ao derivar \textbf{todas} as constantes operacionais de constantes fundamentais da natureza --- a constante de estrutura fina $\afine = 1/137{,}036$ e a raiz da base do logaritmo natural $\sqrt{e} = 1{,}64872$ --- através da fatoração $\betatgl = \afine \times \sqrt{e} = 0{,}012031$, denominada Constante de Miguel na Teoria da Gravitação Luminodinâmica.

Adicionalmente, a presente invenção introduz uma etapa de pré-processamento não-linear inédita --- a radicalização luminodinâmica $g = \sqrt{|L|}$ combinada com nominação de estados psiônicos --- que transforma os dados de entrada antes da mistura ARX, proporcionando propriedades de avalanche superiores à ingestão linear de blocos.

% ────────────────────────────────────────────────────────────────────────
\section{Problema Técnico}
% ────────────────────────────────────────────────────────────────────────

O problema técnico resolvido pela presente invenção é a geração de resumos criptográficos de 256 bits com:

\begin{enumerate}[nosep]
\item Zero parâmetros arbitrários --- todas as constantes operacionais derivadas de constantes fundamentais da física;
\item Pré-processamento não-linear dos dados de entrada via operação radical ($\sqrt{|\cdot|}$) e classificação estrutural (nominação psiônica);
\item Propriedade de avalanche estatisticamente ideal (razão média de bits alterados $\approx 1{,}0$ para o ideal de $H/2$);
\item Resistência a colisões verificável por anti-tautologia (a operação radical não é identidade);
\item Compatibilidade com protocolos de rede holográfica e sistemas de consenso baseados na TGL.
\end{enumerate}

% ────────────────────────────────────────────────────────────────────────
\section{Descrição Detalhada da Invenção}
% ────────────────────────────────────────────────────────────────────────

\subsection{Constantes Fundamentais}

Todas as constantes operacionais do hash são derivadas exclusivamente de:

\begin{align}
\afine &= 7{,}2973525693 \times 10^{-3} \quad \text{(constante de estrutura fina, CODATA 2018)} \\
e &= 2{,}71828\ldots \quad \text{(número de Euler)} \\
\betatgl &= \afine \times \sqrt{e} = 0{,}012031\ldots \quad \text{(Constante de Miguel)}
\end{align}

A partir destas, derivam-se:

\begin{itemize}[nosep]
\item \textbf{Número de rounds de mistura por bloco:} $N_R = \lceil\sqrt{1/\betatgl}\rceil = 10$
\item \textbf{Número de rounds de finalização:} $N_F = 4$
\item \textbf{Constantes de rotação (domínio $\afine$):} $r_i = \lfloor \afine \cdot (i+1) \cdot 1000 \rfloor \bmod 63 + 1$, para $i = 0, \ldots, N_R + N_F - 1$
\item \textbf{Constantes multiplicativas (domínio $\sqrt{e}$):} Representação em inteiro de 64 bits de $\sqrt{e} \times 2^{i+10}$, com bit menos significativo forçado a 1 (garantindo inversibilidade modular)
\item \textbf{Chaves de round:} Representações em inteiro de 64 bits de funções alternadas de $\afine$ e $\sqrt{e}$
\item \textbf{Estado inicial:} 32 inteiros de 64 bits derivados de potências cruzadas de $\afine$ e $\sqrt{e}$
\end{itemize}

A fatoração $\betatgl = \afine \times \sqrt{e}$ separa as operações em dois domínios:
\begin{itemize}[nosep]
\item \textbf{Domínio $\afine$ (eletromagnético):} Governa rotações --- difusão posicional dos bits
\item \textbf{Domínio $\sqrt{e}$ (entrópico):} Governa multiplicações --- confusão de valores
\end{itemize}

\subsection{Pipeline de Processamento}

O método compreende quatro etapas sequenciais:

\subsubsection{Etapa 1 --- Preparação (Merkle-Damgård)}

\begin{enumerate}[nosep]
\item Receber dados de entrada de tamanho arbitrário (incluindo vazio)
\item Acrescentar byte de terminação (0x80)
\item Preencher com padrão derivado de $\betatgl$ (representação em bytes do número $\betatgl$ em formato IEEE 754 de 64 bits, ciclando)
\item Acrescentar comprimento original em 8 bytes (\textit{little-endian})
\item Dividir em blocos de 256 bytes (32 inteiros de 64 bits)
\end{enumerate}

\subsubsection{Etapa 2 --- Ingestão com Perturbação TGL}

Para cada bloco de 256 bytes:

\begin{enumerate}[nosep]
\item Leitura direta como 32 inteiros de 64 bits (\textit{little-endian}), preservando toda informação dos bytes originais
\item Interpretação simultânea como 32 valores de ponto flutuante de 64 bits para extração de perturbação TGL:
\begin{enumerate}[nosep]
\item Radicalização: $g_i = \sqrt{|L_i| + \varepsilon}$ (equação primordial da TGL)
\item Sinais: $s_i = \mathrm{sign}(L_i)$ e derivada local $\partial L_i = L_i - L_{i-1}$
\item Nominação psiônica: $\psi_i = f(\mathrm{sign}(L_i), \mathrm{sign}(\partial L_i)) \in \{0,1,2,3\} \subset \mathcal{H}_4$
\item Força de Expulsão: $F_{\mathrm{exp}} = \langle \mathrm{paridade}(\psi) \rangle$
\end{enumerate}
\item Aplicação da perturbação a cada inteiro de 64 bits:
\begin{enumerate}[nosep]
\item Rotação à esquerda por $(\psi_i \times 16 + f_{\mathrm{shift}})$ posições
\item XOR com representação inteira de $g_i$
\item Adição modular com constante de perturbação TGL pré-computada
\end{enumerate}
\end{enumerate}

O dado original (uint64) é preservado e apenas \textbf{perturbado} pela informação TGL (radicalização + nominação), garantindo que diferenças mínimas no input se propaguem.

\subsubsection{Etapa 3 --- Mistura ARX ($N_R$ rounds por bloco)}

Para cada round $r = 0, \ldots, N_R - 1$:

\begin{enumerate}[nosep]
\item \textbf{Injeção + Multiplicação:} Para cada palavra $s_i$ do estado: $s_i \leftarrow (s_i \oplus w_i) \times m_r \pmod{2^{64}}$, onde $w_i$ é a palavra do bloco ingerido e $m_r$ é a constante multiplicativa derivada de $\sqrt{e}$ para o round $r$

\item \textbf{Quarter-rounds em colunas:} Para $c = 0, \ldots, 7$, aplicar quarter-round ARX sobre $(s_c, s_{c+8}, s_{c+16}, s_{c+24})$ com rotações derivadas de $\afine$:
\begin{align}
a &\leftarrow a + b; \quad d \leftarrow d \oplus a; \quad d \leftarrow d \lll r_1 \nonumber \\
c &\leftarrow c + d; \quad b \leftarrow b \oplus c; \quad b \leftarrow b \lll r_2 \nonumber \\
a &\leftarrow a + b; \quad d \leftarrow d \oplus a; \quad d \leftarrow d \ggg (r_1/2 + 1) \nonumber \\
c &\leftarrow c + d; \quad b \leftarrow b \oplus c; \quad b \leftarrow b \ggg (r_2/2 + 1) \nonumber
\end{align}

\item \textbf{Quarter-rounds em diagonais:} Para $d = 0, \ldots, 7$, aplicar quarter-round ARX sobre $(s_d, s_{(d+1)\bmod 8 + 8}, s_{(d+2)\bmod 8 + 16}, s_{(d+3)\bmod 8 + 24})$

\item \textbf{Adição de chave de round:} $s_i \leftarrow s_i + K_{(i + r) \bmod 32} \pmod{2^{64}}$
\end{enumerate}

\subsubsection{Etapa 4 --- Finalização}

Após processar todos os blocos:

\begin{enumerate}[nosep]
\item \textbf{Difusão pura ($N_F = 4$ rounds):} Mesma estrutura ARX da Etapa 3, mas sem injeção de bloco. Em vez disso, auto-mistura: $s_i \leftarrow s_i \oplus s_{(i+7) \bmod 32}$ antes de cada round.

\item \textbf{XOR-fold:} Redução do estado de 32 palavras para 4 palavras:
\begin{align}
f_i^{(8)} &= s_i \oplus s_{i+8} \oplus s_{i+16} \oplus s_{i+24}, \quad i = 0,\ldots,7 \nonumber \\
f_i^{(4)} &= f_i^{(8)} \oplus f_{i+4}^{(8)}, \quad i = 0,\ldots,3 \nonumber
\end{align}

\item \textbf{Serialização:} Concatenação de $f_0^{(4)}, f_1^{(4)}, f_2^{(4)}, f_3^{(4)}$ em \textit{little-endian}, produzindo 32 bytes (256 bits).
\end{enumerate}

% ────────────────────────────────────────────────────────────────────────
\section{Resultados Experimentais}
% ────────────────────────────────────────────────────────────────────────

A implementação de referência foi testada com suite de validação de 7 testes:

\begin{table}[H]
\centering
\caption{Resultados do benchmark TGL-Hash v1.2.}
\label{tab:resultados}
\begin{tabular}{lccl}
\toprule
\textbf{Teste} & \textbf{Resultado} & \textbf{Critério} & \textbf{Status} \\
\midrule
Determinismo & 100/100 & Todos iguais & APROVADO \\
Avalanche & 128,1/128 (razão 1,001) & $0{,}4 < r < 1{,}6$ & APROVADO \\
Distribuição $\chi^2$ & 245,5 (32/32 posições) & $\chi^2 < 293$ & APROVADO \\
Colisões & 0/50.000 & Zero colisões & APROVADO \\
Sensibilidade & 6/6 pares distintos & Todos diferentes & APROVADO \\
Velocidade (1 MB, CPU) & 1,03 MB/s & $> 0{,}5$ MB/s & APROVADO \\
Inputs especiais & 5/5 & Todos corretos & APROVADO \\
\bottomrule
\end{tabular}
\end{table}

A propriedade de avalanche (razão 1,001 com desvio padrão 7,97) é estatisticamente indistinguível do ideal teórico, demonstrando que a combinação de perturbação TGL e mistura ARX com constantes fatoradas produz difusão completa.

% ────────────────────────────────────────────────────────────────────────
\section{Vantagens da Invenção}
% ────────────────────────────────────────────────────────────────────────

\begin{enumerate}[nosep]
\item \textbf{Zero parâmetros arbitrários:} Todas as constantes operacionais são derivadas de $\afine$ e $\sqrt{e}$, constantes fundamentais da natureza com valores conhecidos com precisão de $10^{-10}$.
\item \textbf{Pré-processamento não-linear:} A radicalização $g = \sqrt{|L|}$ e a nominação psiônica $\psi \in \mathcal{H}_4$ transformam os dados antes da mistura, proporcionando fonte adicional de não-linearidade ausente em funções hash convencionais.
\item \textbf{Separação de domínios:} A fatoração $\betatgl = \afine \times \sqrt{e}$ separa algebricamente as operações de difusão ($\afine$, rotações) e confusão ($\sqrt{e}$, multiplicações), permitindo análise independente de cada componente.
\item \textbf{Avalanche estatisticamente perfeita:} Razão medida de 1,001 (ideal 1,000), superior à maioria das implementações de SHA-256 em dados estruturados.
\item \textbf{Compatibilidade com protocolos TGL:} Integrável ao protocolo de rede $\Psi$-NET, ao sistema de consenso TGL-Coin, e ao sistema de integridade ACOM Crypto.
\end{enumerate}

% ────────────────────────────────────────────────────────────────────────
\section{Aplicações Industriais}
% ────────────────────────────────────────────────────────────────────────

\begin{itemize}[nosep]
\item Integridade de pacotes em protocolos de rede holográfica ($\Psi$-NET)
\item Assinatura de blocos em sistemas de consenso distribuído (TGL-Coin)
\item Verificação de integridade em sistemas de compressão (ACOM)
\item Validação de separação dimensional em criptografia ontológica (ACOM Crypto)
\item Geração de identificadores únicos para dados científicos
\item Detecção de alterações em modelos de inteligência artificial (checksums de pesos)
\end{itemize}

% ────────────────────────────────────────────────────────────────────────
\section{Referências à Arte Prévia}
% ────────────────────────────────────────────────────────────────────────

\begin{enumerate}[nosep]
\item NIST FIPS 180-4: Secure Hash Standard (SHA-256), 2015.
\item NIST FIPS 202: SHA-3 Standard (Keccak), 2015.
\item O'Connor, J. et al. BLAKE3, 2020.
\item Bernstein, D.J. The ChaCha family of stream ciphers, 2008.
\item Miguel, L.A.R. A Fatoração da Constante de Miguel: $\betatgl = \afine \times \sqrt{e}$. Zenodo, 2026. DOI: 10.5281/zenodo.18852146.
\item Patente BR 10 2026 003431 2 --- Método e Sistema de Reflexão Dimensional de Dados (ACOM Mirror/Crypto). INPI, 2026.
\item Patente BR 10 2026 003453 3 --- Protocolo de Internet Luminodinâmico ($\Psi$-NET). INPI, 2026.
\end{enumerate}

% ============================================================================
% REIVINDICAÇÕES
% ============================================================================

\newpage
\section*{Quadro Reivindicatório}

\begin{enumerate}
\item \textbf{MÉTODO DE GERAÇÃO DE RESUMO CRIPTOGRÁFICO (HASH) DE 256 BITS} compreendendo as etapas de: receber dados de entrada de tamanho arbitrário; preparar os dados com preenchimento derivado da Constante de Miguel $\betatgl = \afine \times \sqrt{e}$ e acréscimo de comprimento; dividir em blocos de tamanho fixo; para cada bloco, ler como inteiros de 64 bits e aplicar perturbação não-linear compreendendo radicalização $g = \sqrt{|L|}$ e nominação de estados psiônicos $\psi \in \mathcal{H}_4$; processar cada bloco com $N_R = \lceil\sqrt{1/\betatgl}\rceil$ rounds de mistura ARX (\textit{Add-Rotate-XOR}) com constantes de rotação derivadas de $\afine$ e constantes multiplicativas derivadas de $\sqrt{e}$; aplicar $N_F$ rounds de difusão pura; e extrair o resumo de 256 bits por XOR-fold do estado interno; o método caracterizado por todas as constantes operacionais serem derivadas exclusivamente das constantes de estrutura fina $\afine$ e do número de Euler $e$ através da fatoração $\betatgl = \afine \times \sqrt{e}$.

\item Método, de acordo com a reivindicação 1, caracterizado pelo fato de a perturbação TGL compreender: interpretação simultânea dos bytes do bloco como valores de ponto flutuante de 64 bits, extração da componente gravitacional $g_i = \sqrt{|L_i|}$, determinação de sinais $s_i = \mathrm{sign}(L_i)$ e derivada local $\partial L_i$, e nominação de estado psiônico $\psi_i = f(s_i, \mathrm{sign}(\partial L_i)) \in \{0,1,2,3\}$ no Espaço de Hilbert tetradimensional $\mathcal{H}_4$.

\item Método, de acordo com a reivindicação 1, caracterizado pelo fato de o número de rounds de mistura por bloco ser $N_R = \lceil\sqrt{1/\betatgl}\rceil$, onde $\betatgl = \afine \times \sqrt{e}$, resultando em $N_R = 10$ rounds derivados da amplificação holográfica.

\item Método, de acordo com a reivindicação 1, caracterizado pelo fato de as constantes de rotação serem $r_i = \lfloor \afine \cdot (i+1) \cdot 1000 \rfloor \bmod 63 + 1$, derivadas do domínio eletromagnético (constante de estrutura fina), e as constantes multiplicativas serem derivadas da representação inteira de $\sqrt{e} \times 2^{i+10}$ com bit menos significativo forçado a 1, derivadas do domínio entrópico.

\item Método, de acordo com a reivindicação 1, caracterizado pelo fato de cada round de mistura compreender: injeção por XOR do bloco ingerido, multiplicação modular por constante derivada de $\sqrt{e}$, quarter-rounds ARX em colunas sobre grupos de 4 palavras de 64 bits $(s_c, s_{c+8}, s_{c+16}, s_{c+24})$, quarter-rounds ARX em diagonais, e adição de chave de round derivada de $\afine$ e $\sqrt{e}$ alternados.

\item Método, de acordo com a reivindicação 1, caracterizado pelo fato de a etapa de finalização compreender $N_F = 4$ rounds de difusão pura (sem injeção de bloco) com auto-mistura $s_i \leftarrow s_i \oplus s_{(i+7) \bmod 32}$, seguida de XOR-fold que reduz o estado de 32 palavras para 4 palavras de 64 bits.

\item Método, de acordo com a reivindicação 1, caracterizado pelo fato de o preenchimento utilizar padrão derivado da representação em formato IEEE 754 de 64 bits da constante $\betatgl$, em conformidade com a construção Merkle-Damgård.

\item Método, de acordo com a reivindicação 1, caracterizado pelo fato de suportar operação incremental (\textit{streaming}), onde dados podem ser acrescentados progressivamente antes da extração do resumo final.

\item Método, de acordo com a reivindicação 2, caracterizado pelo fato de a Força de Expulsão $F_{\mathrm{exp}}$ do bloco, derivada da distribuição de paridades dos estados psiônicos, ser utilizada como parâmetro de rotação adicional na perturbação, introduzindo dependência não-linear entre a estrutura estatística do bloco e a transformação aplicada.

\item \textbf{SISTEMA DE GERAÇÃO DE RESUMO CRIPTOGRÁFICO} compreendendo: um módulo de preparação configurado para particionar dados de entrada em blocos com preenchimento derivado de $\betatgl$; um módulo de ingestão configurado para converter blocos em inteiros de 64 bits e aplicar perturbação via radicalização $g = \sqrt{|L|}$ e nominação psiônica $\psi \in \mathcal{H}_4$; um módulo de mistura ARX configurado para executar $N_R = \lceil\sqrt{1/\betatgl}\rceil$ rounds por bloco com constantes derivadas de $\afine$ (rotações) e $\sqrt{e}$ (multiplicações); um módulo de finalização configurado para executar $N_F$ rounds de difusão pura e extrair 256 bits por XOR-fold; o sistema caracterizado por todas as constantes operacionais serem derivadas exclusivamente de $\afine$ e $e$ através da fatoração $\betatgl = \afine \times \sqrt{e}$, resultando em zero parâmetros arbitrários.

\item Sistema, de acordo com a reivindicação 10, caracterizado por suportar execução em CPU via operações inteiras de 64 bits e em GPU via operações tensoriais, produzindo resultados idênticos.

\item Sistema, de acordo com a reivindicação 10, caracterizado por ser integrável como função de integridade no protocolo de rede $\Psi$-NET, no sistema de consenso TGL-Coin, e no sistema de verificação ACOM Crypto.
\end{enumerate}

% ============================================================================
% RESUMO
% ============================================================================

\newpage
\section*{Resumo}

\textbf{MÉTODO E SISTEMA DE GERAÇÃO DE RESUMO CRIPTOGRÁFICO (HASH) BASEADO EM RADICALIZAÇÃO LUMINODINÂMICA, NOMINAÇÃO PSIÔNICA E MISTURA ARX COM CONSTANTES FUNDAMENTAIS FATORADAS}

\medskip

A presente invenção refere-se a métodos e sistemas para geração de resumos criptográficos de 256 bits, aplicáveis aos setores de segurança da informação, integridade de dados e protocolos de rede. O método compreende receber dados de entrada de tamanho arbitrário, aplicar perturbação não-linear via radicalização $g = \sqrt{|L|}$ e nominação de estados psiônicos $\psi$ no Espaço de Hilbert tetradimensional $\mathcal{H}_4$, processar com $N_R = \lceil\sqrt{1/\betatgl}\rceil = 10$ rounds de mistura ARX (\textit{Add-Rotate-XOR}) com constantes de rotação derivadas da constante de estrutura fina $\afine = 1/137{,}036$ e constantes multiplicativas derivadas de $\sqrt{e} = 1{,}64872$, aplicar 4 rounds de difusão pura, e extrair o resumo por XOR-fold. Todas as constantes operacionais são derivadas exclusivamente da fatoração $\betatgl = \afine \times \sqrt{e} = 0{,}012031$, resultando em zero parâmetros arbitrários. Testes demonstram avalanche com razão 1,001 (ideal 1,000), distribuição $\chi^2$ aprovada em 32/32 posições de byte, zero colisões em 50.000 hashes, e throughput de 1,03~MB/s em CPU.

\medskip
\noindent Figura principal: Figura 1 (Pipeline do Método)

\end{document}
